\subsection{BOUNDS FOR THE REMAINDER IN THE EULER-MACLAURIN SUMMATION FORMULA}

The estimate obtained in (16) for the Bernoulli polynomials on the interval \([0, 1]\) allows one to estimate the remainder \(T_{p}(x, n)\) in the Euler-Maclaurin summation formula (VI, p.~282, formula (39)) easily:

\[
\left\{ \begin{array}{l} f (x) + f (x + 1) + \dots + f (x + n) \\ = \int_ {x} ^ {x + n + 1} f (t) d t - \frac {1}{2} \bigl (f (x + n + 1) - f (x) \bigr) \\ \qquad + \sum_ {k = 1} ^ {p} \frac {b _ {2 k}}{(2 k) !} \bigl (f ^ {(2 k - 1)} (x + n + 1) - f ^ {(2 k - 1)} (x) \bigr) + T _ {p} (x, n). \end{array} \right.\tag{1}
\]

Indeed, one has (VI, p.~283, formula (41))

\[
\mathrm{T} _ {p} (x, n) = - \frac {1}{(2 p + 1) !} \int_ {0} ^ {n + 1} \overline {{\mathrm{B}}} _ {2 p + 1} (1 - s) f ^ {(2 p + 1)} (x + s) d s\tag{2}
\]

where \(\overline{\mathbf{B}}_{2p + 1}(t)\) is the periodic function of period 1 equal to \(\mathrm{B}_{2p + 1}(t)\) on the interval {[}0, 1{[}. The formula (16) of VI, p.~288, shows that

\[
\left| \overline {{{{\mathrm{B}}}}} _ {2 p + 1} (t) \right| \leqslant 4 e ^ {2 \pi} \frac {(2 p + 1) !}{(2 \pi) ^ {2 p + 1}}\tag{3}
\]

for all \(t \in R\) , and applying the mean value formula gives the estimate

\[
\left| \mathrm{T} _ {p} (x, n) \right| \leqslant \frac {4 e ^ {2 \pi}}{(2 \pi) ^ {2 p + 1}} \int_ {x} ^ {x + n + 1} \left| f ^ {(2 p + 1)} (t) \right| d t\tag{4}
\]

for \(\mathrm{T}_p(x,n)\) .

